\documentclass[12pt,letterpaper]{article}

\usepackage[english]{babel}
\usepackage[utf8]{inputenc}
\usepackage{amsmath}
\usepackage{amsfonts}
\usepackage{amssymb}
\usepackage{graphicx}
\usepackage[colorinlistoftodos]{todonotes}
\usepackage[margin=1in]{geometry}
\usepackage{wrapfig}
\usepackage[justification=centering]{caption}
\usepackage[subnum]{cases}
\setlength{\parindent}{0pt}
\setlength{\parskip}{1em}
\usepackage{empheq}


\title{Agron 2002 Statistics\\ {\bf Lab Report 2}}
\date{23 March 2026}


\begin{document}
\maketitle
\section*{Instruction}
\begin{itemize}
    \item Paste your code (text, not screenshot) and output (screenshot) in a pdf file
    \item Format: \textbf{pdf} only
    \item Deadline: 2026/03/29 (SUN) 23:59
    \item Late submission is not accepted
\end{itemize}


\section*{Exercise 1: Descriptive Statistics}
Your friend is planning to buy a new mobile phone as a gift for the person he/she is dating. 
You found a file named \textit{\textbf{mobiles.csv}} on the internet that contains details about various mobile 
phones. Explore the data to get useful information:

\textbf{a.} How many mobile phone models are included in the dataset?

\textbf{b.} If your friend wants the screen size to be greater than the 42th percentile but less than the 87th percentile, what range of screen sizes should you look for?

\textbf{c.} Your friend wants to know about the descriptive statistics of mobile phone weights. Use 
the functions summary() and var() to get statistics for the "Mobile Weight" column. 
What do you get?

\textbf{d.} Use the boxplot() function to visualize the data for the "Mobile Weight" column. What 
do you observe? Do the results align with those from (c)?

\textbf{e.} In which currency can you buy the cheapest mobile phone, and how much does it cost?

Note:\\
*The data is on NTU COOL. Use read\_csv() to read this data and call it \textit{\textbf{mymobiles}}. \\
*Units: Screen size: inches; Phone weights: gram \\
*1 USD (USA) = 3.67 AED (Dubai) = 87.33 INR (India)



\section*{Exercise 2: Data Visualization}
The \textit{\textbf{Anderson.csv}} dataset includes 42 leukemia patients, focusing on the time to relapse after receiving treatment. \\
Sex: 0 = female, 1 = male\\
Rx: 0 = New treatment, 1 = Standard treatment\\
Relapse: 0 = No relapse, 1 = Relapse occurred\\
Survt: Time in weeks from the start of treatment until the relapse or the last observation\\
logWBC: white blood cell count (log transform) \\

Visualize the distribution of the data. The figure should be easy to read.\\

\textbf{a.} Plot the bar plot of Sex in the Anderson dataset. Set the figure title to “Bar Plot of Sex.” \\
Modify the y-axis limit to range from 0 to 25. Change the color of the bars to whatever you like.

\textbf{b.} Plot the histogram of logWBC. Change the number of bins to whatever you like.

\textbf{c.} Plot the scatter plot of logWBC (x) vs. Survt (y). Set the x-axis label to “logWBC” and the y-axis label to “Survival Time.”

Note: \\
The data can be found on NTU COOL. Use read\_csv() to read this data and call it \textit{\textbf{anderson}}



\end{document}